\documentclass[11pt]{article}

\usepackage[a4paper,margin=1in]{geometry}
\usepackage{amsmath,amssymb}
\usepackage[T1]{fontenc}
\usepackage{lmodern}
\usepackage{hyperref}

\title{Methods and Assumptions for the Figure 9 Sensitivity Curves and SNR Table}
\author{}
\date{\today}

\begin{document}
\maketitle

\section{Scope}
This note documents the equations, assumptions, and numerical settings used in the \texttt{felix} Figure~9 workflow:
\begin{itemize}
\item the sensitivity-curve notebook \texttt{FIGUUR9\_[JUISTE\_VOOR\_BERT].ipynb},
\item the equivalent script \texttt{notebook\_code.py}, which contains the same Figure~9 curve logic and also the SNR-table calculation.
\end{itemize}

The notebook itself only plots the sensitivity curves. The SNR table is implemented in \texttt{notebook\_code.py}. The formulas are the same in both places; the only relevant run-time difference is that the notebook uses \texttt{nside=800} while \texttt{notebook\_code.py} uses \texttt{NSIDE\_FAST=820}.

\section{Exact Numerical Setup}
\subsection{Common constants}
The code fixes
\begin{align}
c &= 299\,792\,458~\mathrm{m\,s^{-1}}, \\
L &= 2.5\times 10^9~\mathrm{m}, \\
f_* &= \frac{c}{2\pi L}, \\
M_{\rm pc} &= 3.0856775814913673\times 10^{22}~\mathrm{m}, \\
\frac{H_0}{h} &= \frac{100\times 1000}{M_{\rm pc}}~\mathrm{s^{-1}}, \\
Y_{00} &= \frac{1}{\sqrt{4\pi}}.
\end{align}

\subsection{Figure notebook run}
In \texttt{FIGUUR9\_[JUISTE\_VOOR\_BERT].ipynb} the plotted curves use
\begin{align}
f &= \text{\texttt{np.logspace(-4.5,-0.35,320)}}~\mathrm{Hz}, \\
\ell_{\max} &= 10, \\
\text{\texttt{nside}} &= 800, \\
\text{\texttt{iter\_sht}} &= 0.
\end{align}

\subsection{Equivalent script run}
In \texttt{notebook\_code.py}, the example block uses
\begin{align}
f &= \text{\texttt{np.logspace(-4.5,-0.35,320)}}~\mathrm{Hz}, \\
\ell_{\max} &= 10, \\
\text{\texttt{NSIDE\_FAST}} &= 820, \\
\text{\texttt{ITER\_SHT}} &= 0, \\
T_{\rm obs} &= 4~\mathrm{yr}.
\end{align}

For the SNR table in \texttt{notebook\_code.py}, the signal model is taken to be
\begin{align}
\Omega_{\rm GW}(f) h^2 &= 10^{-12}\left(\frac{f}{2.5\times 10^{-3}~\mathrm{Hz}}\right)^0 = 10^{-12}, \\
C_\ell^{\rm GW} &= 1 \qquad \text{for all } \ell=0,\dots,10.
\end{align}

\section{Geometric Assumptions}
The LISA geometry used here is \emph{static}. There is no orbital motion, no time averaging, and no annual modulation. The three spacecraft positions are fixed to
\begin{align}
\mathbf{r}_1 &= (0,0,0), \\
\mathbf{r}_2 &= (L,0,0), \\
\mathbf{r}_3 &= \left(\frac{L}{2},\frac{\sqrt{3}L}{2},0\right).
\end{align}

The arm directions are therefore
\begin{align}
\hat{\ell}_{12} &= \frac{\mathbf{r}_2-\mathbf{r}_1}{|\mathbf{r}_2-\mathbf{r}_1|}, &
\hat{\ell}_{13} &= \frac{\mathbf{r}_3-\mathbf{r}_1}{|\mathbf{r}_3-\mathbf{r}_1|}, \\
\hat{\ell}_{23} &= \frac{\mathbf{r}_3-\mathbf{r}_2}{|\mathbf{r}_3-\mathbf{r}_2|}, &
\hat{\ell}_{21} &= \frac{\mathbf{r}_1-\mathbf{r}_2}{|\mathbf{r}_1-\mathbf{r}_2|}, \\
\hat{\ell}_{31} &= \frac{\mathbf{r}_1-\mathbf{r}_3}{|\mathbf{r}_1-\mathbf{r}_3|}, &
\hat{\ell}_{32} &= \frac{\mathbf{r}_2-\mathbf{r}_3}{|\mathbf{r}_2-\mathbf{r}_3|}.
\end{align}

Each spacecraft uses the two outgoing/incoming arms
\begin{align}
1 &: (\hat{\ell}_{12},\hat{\ell}_{13}), \\
2 &: (\hat{\ell}_{23},\hat{\ell}_{21}), \\
3 &: (\hat{\ell}_{31},\hat{\ell}_{32}).
\end{align}

The sky is discretized on a HEALPix grid with \texttt{nest=False}, i.e.\ ring ordering.

\section{Noise Model}
The code uses the Appendix~B noise functions
\begin{align}
S_{\rm oms}(f) &=
\left(15\times 10^{-12}\right)^2
\left[1+\left(\frac{2\times 10^{-3}}{f}\right)^4\right], \\
S_{\rm acc}(f) &=
\left(3\times 10^{-15}\right)^2
\left[1+\left(\frac{0.4\times 10^{-3}}{f}\right)^2\right]
\left[1+\left(\frac{f}{8\times 10^{-3}}\right)^4\right].
\end{align}

From these, the channel noise levels are
\begin{align}
\widetilde N_{AE}(f) &=
\frac{1}{2}\left(2+\cos\frac{f}{f_*}\right)\frac{S_{\rm oms}(f)}{L^2}
{}+
2\left(1+\cos\frac{f}{f_*}+\cos^2\frac{f}{f_*}\right)
\frac{S_{\rm acc}(f)}{L^2}\left(\frac{1}{2\pi f}\right)^4, \\
\widetilde N_T(f) &=
\left(1-\cos\frac{f}{f_*}\right)\frac{S_{\rm oms}(f)}{L^2}
{}+
2\left(1-\cos\frac{f}{f_*}\right)^2
\frac{S_{\rm acc}(f)}{L^2}\left(\frac{1}{2\pi f}\right)^4.
\end{align}

The code then sets
\begin{equation}
\widetilde N_A(f)=\widetilde N_E(f)=\widetilde N_{AE}(f),
\qquad
\widetilde N_T(f)=\widetilde N_T(f).
\end{equation}

\section{Sky Basis, Polarizations, and Arm Transfer Functions}
For each HEALPix direction $(\theta,\phi)$, the unit wavevector is
\begin{equation}
\hat{\mathbf{k}} =
(\sin\theta\cos\phi,\sin\theta\sin\phi,\cos\theta).
\end{equation}

The basis vectors used in the code are
\begin{align}
\hat{\mathbf{e}}_\theta &=
(\cos\theta\cos\phi,\cos\theta\sin\phi,-\sin\theta), \\
\hat{\mathbf{e}}_\phi &=
(-\sin\phi,\cos\phi,0).
\end{align}

The polarization tensors are
\begin{align}
e^+_{ij} &= (\hat e_\theta)_i(\hat e_\theta)_j-(\hat e_\phi)_i(\hat e_\phi)_j, \\
e^\times_{ij} &= (\hat e_\theta)_i(\hat e_\phi)_j+(\hat e_\phi)_i(\hat e_\theta)_j.
\end{align}

For any arm direction $\hat{\ell}$, the geometric projection used in the code is
\begin{equation}
g^P(\hat{\ell}) = \frac{1}{2}\,\ell_i \ell_j e^P_{ij},
\qquad
P\in\{+,\times\}.
\end{equation}

The transfer-function conventions are
\begin{align}
u &= \frac{f}{f_*}, \\
\mu &= \hat{\mathbf{k}}\cdot \hat{\ell}, \\
\operatorname{sinc}(x) &= \frac{\sin x}{x},
\end{align}
with the explicit code limit $\operatorname{sinc}(0)=1$.

The intermediate functions are
\begin{align}
M(u,\mu) &= e^{\frac{i}{2}u(1+\mu)} \, \operatorname{sinc}\!\left[\frac{u}{2}(1+\mu)\right], \\
T(u,\mu) &= e^{-iu} M(u,-\mu) + e^{-iu\mu} M(u,\mu).
\end{align}

For spacecraft $i$, with two arms $a$ and $b$, the code constructs
\begin{align}
R_i^+(f,\hat{\mathbf{k}}) &= g_a^+\,T(u,\mu_a)-g_b^+\,T(u,\mu_b), \\
R_i^\times(f,\hat{\mathbf{k}}) &= g_a^\times\,T(u,\mu_a)-g_b^\times\,T(u,\mu_b).
\end{align}

\section{A/E/T Mixing}
The Figure~9 code works with the constant mixing matrix
\begin{equation}
C_{\rm AET} =
\begin{pmatrix}
-1/\sqrt{2} & 0 & 1/\sqrt{2} \\
1/\sqrt{6} & -2/\sqrt{6} & 1/\sqrt{6} \\
1/\sqrt{3} & 1/\sqrt{3} & 1/\sqrt{3}
\end{pmatrix},
\end{equation}
whose rows correspond to $(A,E,T)$ and whose columns correspond to the three spacecraft responses.

For a channel pair $(O,O')$, the code uses
\begin{equation}
{\rm Mix}_{OO',ij} = C_{Oi}\,C_{O'j}.
\end{equation}

\section{Response Maps and Spherical Harmonics}
The pairwise polarization sum is
\begin{equation}
\sum_A R_i^A R_j^{A*}
=
R_i^+ R_j^{+*} + R_i^\times R_j^{\times *}.
\end{equation}

The sky map corresponding to Eq.~(4.15) is implemented as
\begin{equation}
f_{ij}(f,\hat{\mathbf{k}})
= \frac{1}{8\pi}
\left(R_i^+ R_j^{+*} + R_i^\times R_j^{\times *}\right)
\exp\!\left[-2\pi i f\,\hat{\mathbf{k}}\cdot \Delta\mathbf{x}_{ij}/c\right].
\end{equation}

The channel-pair sky maps are then
\begin{equation}
f_{OO'}(f,\hat{\mathbf{k}})
= \sum_{i,j} C_{Oi} C_{O'j} \, f_{ij}(f,\hat{\mathbf{k}}).
\end{equation}

Because these maps are complex, the code does \emph{not} use a single real-map harmonic transform. Instead it computes
\begin{align}
a_{\ell m} &= \int d\Omega \, f_{OO'}(\hat{\mathbf{k}})\,Y_{\ell m}^*(\hat{\mathbf{k}}), \\
c_{\ell m} &= \int d\Omega \, f_{OO'}^*(\hat{\mathbf{k}})\,Y_{\ell m}^*(\hat{\mathbf{k}}),
\end{align}
by separately transforming the real and imaginary parts with \texttt{healpy.map2alm}.

The rotation-invariant harmonic power used in the code is
\begin{equation}
P_\ell^{OO'} =
|a_{\ell 0}|^2 + \sum_{m=1}^{\ell}
\left(|a_{\ell m}|^2 + |c_{\ell m}|^2\right).
\end{equation}

The response entering the sensitivity curves is then taken to be
\begin{equation}
\widetilde R_{OO'}^\ell(f) = \sqrt{\pi\,P_\ell^{OO'}}.
\end{equation}

This explicit $\sqrt{\pi}$ factor is an implementation choice stated in the code comments: it is included so that the low-frequency $AA$ monopole limit matches $9/20$ rather than $9/(20\sqrt{\pi})$.

To avoid division by zero, the code floors the response:
\begin{equation}
\widetilde R_{OO'}^\ell(f) \ge 10^{-50}
\quad \text{in } \texttt{FIGUUR9\_[JUISTE\_VOOR\_BERT].ipynb},
\end{equation}
and
\begin{equation}
\widetilde R_{OO'}^\ell(f) \ge 10^{-45}
\quad \text{in } \texttt{notebook\_code.py}.
\end{equation}

\section{Sensitivity Curves}
For a given ordered pair $(O,O')$, the code defines
\begin{equation}
\Omega_{{\rm GW},n}^{OO',\ell}(f) h^2
=
\frac{4\pi^2\sqrt{4\pi}}{3(H_0/h)^2}\,
f^3
\frac{\sqrt{\widetilde N_O(f)\widetilde N_{O'}(f)}}{\widetilde R_{OO'}^\ell(f)}.
\end{equation}

The code then forms the inverse-variance combination over \emph{all nine ordered pairs}
\begin{equation}
(O,O') \in \{AA,AE,AT,EA,EE,ET,TA,TE,TT\},
\end{equation}
namely
\begin{equation}
\Omega_{{\rm GW},n}^{\ell,{\rm raw}}(f) h^2
=
\left[
\sum_{O,O'}
\frac{1}{\left(\Omega_{{\rm GW},n}^{OO',\ell}(f) h^2\right)^2}
\right]^{-1/2}.
\end{equation}

This means that the implementation \emph{does} count both $AE$ and $EA$ separately, because it explicitly loops over ordered pairs.

The plotted Figure~9 curve is the rescaled version
\begin{equation}
\Omega_{{\rm GW},n}^{\ell,{\rm Fig9}}(f) h^2
=
Y_{00}\,
\Omega_{{\rm GW},n}^{\ell,{\rm raw}}(f) h^2
=
\frac{1}{\sqrt{4\pi}}\,
\Omega_{{\rm GW},n}^{\ell,{\rm raw}}(f) h^2.
\end{equation}

So the notebook plot label
\begin{equation}
\Omega_{{\rm GW},n}^{\ell}(f)\,h^2/\sqrt{4\pi}
\end{equation}
is literally the quantity that is shown.

\section{SNR Table}
The SNR table is computed in \texttt{notebook\_code.py} through \texttt{compute\_snr\_per\_multipole\_eq444}. The implemented equation is
\begin{equation}
{\rm SNR}_\ell^2
=
T_{\rm obs}
\int df
\left[
\frac{\Omega_{\rm GW}^{\ell}(f) h^2}
\Omega_{{\rm GW},n}^{\ell,{\rm raw}}(f) h^2
\right]^2,
\end{equation}
with
\begin{equation}
\Omega_{\rm GW}^{\ell}(f) h^2
=
\sqrt{C_\ell^{\rm GW}}\,
\Omega_{\rm GW}(f) h^2.
\end{equation}

For the example table in \texttt{notebook\_code.py}, the choices are
\begin{align}
T_{\rm obs} &= 4\times 365.25\times 24\times 3600~\mathrm{s}, \\
\Omega_{\rm GW}(f) h^2 &= 10^{-12}, \\
C_\ell^{\rm GW} &= 1.
\end{align}

The integral is evaluated with the trapezoidal rule on the same frequency grid used for the sensitivity curves:
\begin{equation}
{\rm SNR}_\ell^2
\approx
T_{\rm obs}
\operatorname{trapz}\!\left[
\left(
\frac{\Omega_{\rm GW}^{\ell}(f_k) h^2}
\Omega_{{\rm GW},n}^{\ell,{\rm raw}}(f_k) h^2
\right)^2,
f_k
\right].
\end{equation}

The SNR itself is then
\begin{equation}
{\rm SNR}_\ell = \sqrt{{\rm SNR}_\ell^2}.
\end{equation}

\section{Important Practical Point}
The SNR table must use the \emph{raw} sensitivity curves,
\begin{equation}
\Omega_{{\rm GW},n}^{\ell,{\rm raw}}(f) h^2,
\end{equation}
not the plotted Figure~9 curves divided by $\sqrt{4\pi}$.

Equivalently, if one starts from the plotted curves, one must undo the plotting rescaling:
\begin{equation}
\Omega_{{\rm GW},n}^{\ell,{\rm raw}}(f) h^2
=
\sqrt{4\pi}\,
\Omega_{{\rm GW},n}^{\ell,{\rm Fig9}}(f) h^2.
\end{equation}

This is exactly why \texttt{notebook\_code.py} keeps both
\begin{itemize}
\item \texttt{sensitivity["raw"]}, and
\item \texttt{sensitivity["figure9"]}.
\end{itemize}
The SNR table uses \texttt{raw}; the plot uses \texttt{figure9}.

\section{What Is and Is Not Included}
Included in this workflow:
\begin{itemize}
\item finite-arm transfer functions,
\item both $+$ and $\times$ polarizations,
\item baseline phase factors,
\item all nine ordered A/E/T channel pairs,
\item complex-sky harmonic power with both positive and negative $m$ information retained through the $(a_{\ell m},c_{\ell m})$ construction,
\item a static equilateral LISA triangle.
\end{itemize}

Not included in this workflow:
\begin{itemize}
\item LISA orbital motion,
\item time averaging over the constellation motion,
\item pixel-window deconvolution in the final \texttt{FIGUUR9\_[JUISTE\_VOOR\_BERT].ipynb} pipeline,
\item any sky mask or anisotropic detector motion,
\item any SGWB model more complicated than the user-supplied $\Omega_{\rm GW}(f)$ and $C_\ell^{\rm GW}$.
\end{itemize}

\section{Summary}
The Figure~9 curves are obtained by computing static LISA response maps on a HEALPix sky, projecting them onto A/E/T channel pairs, extracting the multipole response $\widetilde R^\ell_{OO'}(f)$, converting that response into ordered-pair noise curves, inverse-variance combining the nine ordered pairs, and finally dividing by $\sqrt{4\pi}$ for plotting.

The SNR table uses the same pipeline but replaces the plotted curves by the raw, unscaled sensitivity curves and inserts them into Eq.~(4.44) with the specific example choices
\begin{equation}
T_{\rm obs}=4~{\rm yr},
\qquad
\Omega_{\rm GW}(f)h^2=10^{-12},
\qquad
C_\ell^{\rm GW}=1.
\end{equation}

\end{document}
